\section{Executable improvement from public information}
\label{sec:incumbent-tail}
The recovered opponent model supports a direct alternative to learning a
response: evaluate complete endgame policies under the hidden-card distribution
consistent with public play. We implement this approach at the first learner
decision in the final round with at most five draws remaining, provided a
preceding round has already identified the opponent. The retained parent plays
the entire prefix. Figure~\ref{fig:controller-overview} summarizes the
base controller's inference, adoption and execution steps.

\begin{figure}[htbp]
\centering
\begingroup
\footnotesize
\setlength{\fboxsep}{5pt}
\fbox{\begin{minipage}[c][4.5em][c]{0.265\linewidth}
\centering
\textbf{Public evidence}\\[2pt]
Earlier-round opponent identity; current own/public history
\end{minipage}}
\hfill$\longrightarrow$\hfill
\fbox{\begin{minipage}[c][4.5em][c]{0.265\linewidth}
\centering
\textbf{Finite belief}\\[2pt]
History-compatible hands and weighted stock continuations
\end{minipage}}
\hfill$\longrightarrow$\hfill
\fbox{\begin{minipage}[c][4.5em][c]{0.265\linewidth}
\centering
\textbf{Capped search}\\[2pt]
Complete parent reference and retained candidate programs
\end{minipage}}

\smallskip
\makebox[\linewidth][r]{$\big\downarrow$\hspace{0.135\linewidth}}

\smallskip
\fbox{\begin{minipage}{0.91\linewidth}
\centering
\textbf{Exact adoption test}\\[2pt]
Compare completed values under the same belief: $Q(\sigma)>Q(\pi)$.
\end{minipage}}

\smallskip
\begin{tabular}{@{}c@{\hspace{0.04\linewidth}}c@{}}
$\downarrow$ strict improvement & $\downarrow$ otherwise \\[3pt]
\fbox{\begin{minipage}[c][5em][c]{0.425\linewidth}
\centering
\textbf{Execute the retained program}\\[2pt]
Stored action at each represented history; unchanged parent in a checked
guaranteed-win region.
\end{minipage}}
&
\fbox{\begin{minipage}[c][5em][c]{0.425\linewidth}
\centering
\textbf{Keep the parent}\\[2pt]
Includes missing roots, support refusals, unavailable parent values and
no completed strict improvement.
\end{minipage}}
\end{tabular}
\endgroup
\caption{Base controller: from public play to an executable continuation. Search shares learner
actions across indistinguishable worlds. The parent plays the common prefix;
after adoption, execution follows the retained program, including its checked
parent-continuation rule. The guarantee uses the displayed belief and opponent
model; population inference additionally requires a matching source law.}
\label{fig:controller-overview}
\end{figure}


\subsection{From compatible histories to complete policies}
For each legal hidden hand, we replay the observed opponent actions with the
identified deterministic policy. Compatible hands and the remaining stock
permutations define a finite belief under a uniform valid-deal law. Integer
multiplicities represent equivalent unused stock suffixes. Chance branches
expand only when a draw occurs, and terminal branches contribute their weighted
outcomes immediately.

The planner groups worlds by the learner's own and public history before
choosing an action. Indistinguishable worlds therefore share that choice.
It returns a map from these histories to actions for the complete continuation,
along with an exact rational match value. Executing the map preserves the
information constraints used during evaluation.

We first evaluate root actions followed by the unchanged parent, then search
for policies that also optimize later decisions. Each phase receives a separate
cap of 500,000 live nodes and two million native transitions. An unresolved
branch can bound a candidate's value but cannot supply its executable policy.
When a cap interrupts search, previously completed policies remain available.
The controller adopts the best retained policy only if its exact value exceeds
the unchanged parent's value under the same belief. Otherwise it follows the
parent. A root-only control assigns the combined allowance to the first phase.

\subsection{The conditional improvement guarantee}
Let $H$ be the own/public history at the first eligible root, $\pi$ the parent,
and $\sigma_H$ an adopted complete policy. For final match score
$W\in\{0,\tfrac12,1\}$, the adoption test is
\[
  \Delta(H)
  = \mathbb{E}_{\nu}[W(\sigma_H)\mid H]
    -\mathbb{E}_{\nu}[W(\pi)\mid H] > 0.
\]
Here $\nu$ specifies the deal law, fixed opponent and game rules, and the
evaluated belief is its conditional law given $H$. Set $\Delta(H)=0$ whenever
the controller retains the parent, including unavailable roots or an incomplete
parent evaluation. Because the controller and parent share their entire prefix,
the law of total expectation gives
\[
  \mathbb{E}_{\nu}[W(\mathrm{controller})]
  -\mathbb{E}_{\nu}[W(\pi)]
  =\mathbb{E}_{\nu}[\Delta(H)]\geq 0.
\]
This argument requires commitment to the evaluated continuation and the stated
conditional model. It needs neither a globally optimal policy nor a completed
comparison of every action. Its role is to connect a verified local comparison
to expected match performance.

\subsection{Certified completion of endgame programs}
\label{sec:symbolic-tail}
A complete policy need not list every future action. We extend the explicit
history map with a checked rule that calls the unchanged parent in a region
where every legal continuation has the same match outcome. If $M$ is the stored
map and $C_+$ the certified region, the executed policy is
\[
 \sigma(h)=
 \begin{cases}
   M(h), & h\in\operatorname{dom}(M),\\
   \pi(h), & h\notin\operatorname{dom}(M),\ h\in C_+.
 \end{cases}
\]
The first action remains explicit. Every reachable learner decision must satisfy
one of these cases; the executor rejects any other missing entry. The rule is
part of the evaluated program, so compression preserves commitment to a complete
continuation in the conditional improvement argument.

Let $D$ be the learner's point margin. When only the learner's last turn remains
in the final round, the opponent has no scoring turn. The only adverse transfer
is the exhausted-round dealer award of one point, which changes the margin by
two. Thus $D_{\mathrm{final}}\geq D-2$, and $D>2$ defines $C_+$: every legal
continuation wins. Conversely, when only the opponent's last turn remains,
$D_{\mathrm{final}}\leq D+2$; $D<-2$ certifies a loss and leaves no learner
decision to represent. Both conditions use public turn, round and score fields.

For a certified branch with integer belief mass $m$, search contributes $+m$
or $-m$ directly to its signed outcome sum. It shifts the remaining search
threshold by this contribution and recombines the result before recording a
root value. The executable parent rule supplies the actions omitted in a win
region. Independent Python replay then executes those actions and enumerates
all remaining draws through terminal scoring, rather than reusing the
certificate's predicted outcome. Common completed roots must also retain their
exact values and explicit actions; every omitted history is checked against
the public certificate.

A synthetic three-draw fixture illustrates the resulting coverage change.
With a final-round six-point lead, fixed-continuation evaluation falls from
9,034 to 1,258 native transitions, and contingent evaluation from 9,082 to 1,306.
The selected maps shrink from 462 to six entries and from 251 to six entries,
respectively, plus the parent rule. At a cap of 1,258 transitions per phase,
the explicit controller cannot complete its parent reference. The compressed
program completes that reference at $204/275$ and retains a candidate valued
$461/550$, an independently verified gain of $53/550$ at the same first action.

We separately evaluated the opponent-only loss certificate on all 512 saved
prospective sources under their original allowances
(Appendix~\ref{app:incumbent-development}). It reduces physical
planning transitions by 2.0191\% for the controller and 1.2086\% for root-only
search. Parent-value coverage, adopted policies, conditional gains and realized
scores are unchanged. This complete-cohort computational comparison adds no
source matches; its report is \path{docs/CERTIFIED_TAIL_REPLAY_RESULTS.rst}.

Adding the executable own-win rule to that baseline expands coverage under
the same caps (Table~\ref{tab:symbolic-comparison}). Controller parent references
increase from 226 to 231 of 237 eligible roots, and adoptions from 83 to 86.
Root-only parent references increase from 233 to 236, with 56 adoptions in both
versions. Physical planning transitions fall by 8.84\% and 19.11\%, respectively.
All 978 common completed root values and their retained explicit decisions
match. Compression omits 3,217,165 entries across these programs, each checked
against the public certificate; this count is not a count of unique histories.

\begin{table}[htbp]
\centering
\begin{tabular}{llrrr}
\toprule
Arm & Program & Transitions (M) & Parent values & Adoptions \\
\midrule
Controller & Opponent-only certificate & 223.013 & 226 & 83 \\
Controller & Executable compression & 203.294 & 231 & 86 \\
Root-only & Opponent-only certificate & 148.865 & 233 & 56 \\
Root-only & Executable compression & 120.421 & 236 & 56 \\
\bottomrule
\end{tabular}
\caption{Complete-cohort computational comparison on the same 512 saved sources, including 237 eligible roots. Maximum allowances and fresh-cache rules are unchanged. Transitions count native planning decisions; independent verification is separate.}
\label{tab:symbolic-comparison}
\end{table}


The controller's mean conditional gains over its two parents rise to 0.2651 and
0.3290 percentage points, retaining all 256 sources per parent in the denominator.
Every realized win score and final point margin is unchanged from the
opponent-only baseline. Independent Python verification executes 74,652,353
transitions through terminal scoring, separately from planning. Operational
planning times are 7,599 and 5,745 seconds versus 7,076 and 5,494 for the baseline;
these separate concurrent runs do not establish a latency reduction. The result
is greater certified coverage with fewer native planning transitions. It adds
no source matches or fresh strength test; the original actual-win criterion
remains unmet.

Full counts appear in
\path{docs/SYMBOLIC_TAIL_REPLAY_RESULTS.rst}.

